% El espacio de homomorfismos

\begin{frame}
    \frametitle{Sumar homomorfismos}

    \begin{example}
        Sean $f,g\colon\R^2\to\R^2$ con matrices (en la base canónica)
        \[
            A=\begin{pmatrix}1&2\\0&1\end{pmatrix},\qquad
            A'=\begin{pmatrix}0&1\\3&1\end{pmatrix}.
        \]
        La aplicación $u\mapsto f(u)+g(u)$ es
        \[
            u\longmapsto Au+A'u=(A+A')u,
        \]
        el homomorfismo con matriz
        $A+A'=\begin{pmatrix}1&3\\3&2\end{pmatrix}$.
    \end{example}

    \medskip

    Del mismo modo, $u\mapsto\lambda f(u)$ es el homomorfismo con matriz
    $\lambda A$.

    \medskip

    Los homomorfismos de $V$ en $V'$ forman, pues, un espacio vectorial,
    que vamos a describir.
\end{frame}


\begin{frame}
    \frametitle{El espacio vectorial de los homomorfismos}

    \begin{proposition}
        Sean $V$ y $V'$ espacios vectoriales sobre $k$. Para
        $f,g\in\Hom(V,V')$ y $\lambda\in k$, las aplicaciones
        \[
            (f+g)(u)=f(u)+g(u),\qquad(\lambda f)(u)=\lambda f(u)
        \]
        son homomorfismos. Con estas operaciones, $\Hom(V,V')$ es un
        espacio vectorial sobre $k$, cuyo vector cero es la aplicación
        nula.
    \end{proposition}

    \begin{proof}
        Las aplicaciones de $V$ en $V'$ forman un espacio vectorial con
        estas operaciones (como $k^X$ en el Tema~2), así que basta ver
        que $\Hom(V,V')$ es un subespacio. La aplicación nula es un
        homomorfismo, y para $u,u'\in V$ y $\mu,\nu\in k$,
        \begin{align*}
            (f+g)(\mu u+\nu u')&=\mu f(u)+\nu f(u')+\mu g(u)+\nu g(u')
            =\mu(f+g)(u)+\nu(f+g)(u'),\\
            (\lambda f)(\mu u+\nu u')&=\lambda\bigl(\mu f(u)+\nu f(u')\bigr)
            =\mu(\lambda f)(u)+\nu(\lambda f)(u').\qedhere
        \end{align*}
    \end{proof}
\end{frame}


\begin{frame}
    \frametitle{La matriz de una suma}

    \begin{proposition}
        Sean $f,g\in\Hom(V,V')$, $\lambda\in k$, y sean $B$, $B'$ bases de
        $V$ y $V'$. Entonces
        \[
            M_{B,B'}(f+g)=M_{B,B'}(f)+M_{B,B'}(g),
            \qquad
            M_{B,B'}(\lambda f)=\lambda M_{B,B'}(f).
        \]
    \end{proposition}

    \begin{proof}
        La columna $j$ de $M_{B,B'}(f+g)$ es
        \[
            (f+g)(v_j)_{B'}=\bigl(f(v_j)+g(v_j)\bigr)_{B'}
            =f(v_j)_{B'}+g(v_j)_{B'},
        \]
        que es la suma de las columnas $j$ de $M_{B,B'}(f)$ y $M_{B,B'}(g)$.
        La otra igualdad es análoga.
    \end{proof}
\end{frame}


\begin{frame}
    \frametitle{El espacio de homomorfismos es un espacio de matrices}

    \begin{theorem}
        Sean $B$ y $B'$ bases de $V$ y $V'$, con $\dim V=n$ y
        $\dim V'=m$. La aplicación
        \[
            \Phi\colon\Hom(V,V')\longrightarrow M_{m\times n}(k),
            \qquad f\longmapsto M_{B,B'}(f),
        \]
        es un isomorfismo. En particular,
        \[
            \dim\Hom(V,V')=mn=\dim V\cdot\dim V'.
        \]
    \end{theorem}

    \begin{proof}
        \alert{Lineal:} es la proposición anterior.

        \alert{Inyectiva:} si $M_{B,B'}(f)=0$, entonces $f(v_j)_{B'}=0$, luego
        $f(v_j)=0$ para todo $j$. Como $f$ se anula en una base,
        $f=0$.

        \alert{Sobreyectiva:} dada $A=(a_{ij})$, por el teorema de
        determinación existe un homomorfismo $f$ con
        $f(v_j)=a_{1j}v'_1+\cdots+a_{mj}v'_m$. Entonces
        $M_{B,B'}(f)=A$.
    \end{proof}
\end{frame}


\begin{frame}
    \frametitle{Una base de $\Hom(V,V')$}

    Con las bases $\{v_1,\ldots,v_n\}$ de $V$ y $\{v'_1,\ldots,v'_m\}$ de
    $V'$, sea $f_{ij}\in\Hom(V,V')$ el homomorfismo determinado por
    \[
        f_{ij}(v_l)=
        \begin{cases}
            v'_i&\text{si }l=j,\\
            0&\text{si }l\neq j,
        \end{cases}
        \qquad 1\leq i\leq m,\ 1\leq j\leq n.
    \]
    Su matriz es $E_{ij}$, la matriz con un $1$ en la posición $(i,j)$ y
    ceros en las demás.

    \begin{corollary}
        Los $mn$ homomorfismos $f_{ij}$ son una base de $\Hom(V,V')$, y si
        $M_{B,B'}(f)=(a_{ij})$, entonces
        \[
            f=\sum_{i,j}a_{ij}f_{ij}.
        \]
    \end{corollary}

    Es la base canónica de $M_{m\times n}(k)$ transportada por el
    isomorfismo $\Phi$.
\end{frame}


\begin{frame}
    \frametitle{Ejemplo: una base de $\Hom(\R^2,\R^3)$}

    \begin{example}
        Con las bases canónicas, $\dim\Hom(\R^2,\R^3)=6$, y una base es
        $\{f_{ij}\}$ con $f_{ij}(u_1,u_2)=u_je_i$, para $i=1,2,3$ y
        $j=1,2$.
    \end{example}

    Por ejemplo, $f_{21}(u_1,u_2)=(0,u_1,0)$. El homomorfismo
    $f(u_1,u_2)=(u_1+u_2,\,2u_1-u_2,\,3u_2)$ tiene matriz
    $\begin{pmatrix}1&1\\2&-1\\0&3\end{pmatrix}$, luego
    \[
        f=f_{11}+f_{12}+2f_{21}-f_{22}+3f_{32}.
    \]
    En efecto,
    \[
        (u_1,0,0)+(u_2,0,0)+(0,2u_1,0)-(0,u_2,0)+(0,0,3u_2)
        =(u_1+u_2,\,2u_1-u_2,\,3u_2).
    \]

    \medskip

    Más en general: $\Hom(k^n,k^m)=M_{m\times n}(k)$, y
    $\Hom(k,V')\cong V'$ mediante $f\mapsto f(1)$.
\end{frame}


\begin{frame}
    \frametitle{La composición es bilineal}

    \begin{proposition}
        Sean $f,f_1,f_2\in\Hom(V,V')$ y $g,g_1,g_2\in\Hom(V',V'')$, y
        $\lambda\in k$. Entonces
        \begin{align*}
            g\circ(f_1+f_2)&=g\circ f_1+g\circ f_2,
            &(g_1+g_2)\circ f&=g_1\circ f+g_2\circ f,\\
            (\lambda g)\circ f&=\lambda(g\circ f)=g\circ(\lambda f).
        \end{align*}
    \end{proposition}

    \begin{proof}
        Para $u\in V$: $g(f_1(u)+f_2(u))=g(f_1(u))+g(f_2(u))$ porque $g$ es
        lineal; $(g_1+g_2)(f(u))=g_1(f(u))+g_2(f(u))$ por definición de
        la suma; y las igualdades con $\lambda$ se prueban igual.
    \end{proof}

    \medskip

    Son las propiedades del producto de matrices, que recuperamos al
    traducir a matrices: $M_{B,B''}(g\circ f)=M_{B',B''}(g)\,M_{B,B'}(f)$.
\end{frame}


\begin{frame}
    \frametitle{Endomorfismos y matrices cuadradas}

    \begin{proposition}
        Sea $V$ de dimensión $n$ con una base $B$. La aplicación
        \[
            \End(V)\longrightarrow M_n(k),\qquad f\longmapsto M_{B,B}(f),
        \]
        es un isomorfismo de espacios vectoriales que además cumple
        \[
            M_{B,B}(g\circ f)=M_{B,B}(g)\,M_{B,B}(f),\qquad M_{B,B}(\id_V)=I_n.
        \]
        En particular, $\dim\End(V)=n^2$, y $f$ es un automorfismo si y solo
        si $M_{B,B}(f)$ es invertible.
    \end{proposition}

    \medskip

    La suma y la composición en $\End(V)$ se corresponden con la suma y el
    producto en $M_n(k)$. Como el producto de matrices no es conmutativo
    (Tema~1), tampoco lo es la composición de endomorfismos si $n\geq2$:
    por ejemplo, el endomorfismo de $k^2$ con matriz $E_{12}$ y el de
    matriz $E_{21}$ no conmutan, porque $E_{12}E_{21}=E_{11}\neq
    E_{22}=E_{21}E_{12}$.
\end{frame}
